\chapter{Sensitivity Analyses}
\label{ch:appx-sens}

This appendix documents the four sensitivity analyses run on the deployed pipeline in response to the second phase of the Trentini review. Each analysis perturbs one methodological choice of the Italian or Swedish branch of the pipeline, re-derives the resulting partition, and reports the adjusted Rand index against the deployed partition together with relevant fit diagnostics. The interpretive thresholds are stated upfront for each analysis; the verdicts in Table~\ref{tab:appx-sens-summary} are read against those thresholds rather than as absolute statements of correctness. Detail tables for each analysis are available in the project archive as comma-separated value files under \texttt{v9/outputs/sensitivities/}.


\section{Outlier exclusion threshold (Sweden)}
\label{sec:appx-sens-outlier}

The deployed Mahalanobis outlier-exclusion routine for Sweden (\S\ref{sec:ch3-preprocessing}) operates at $\alpha = 0.001$, the conventional $\chi^2$ cutoff at the 99.9th percentile. The procedure removes 44 respondents from the post-listwise Swedish sub-sample of 1{,}831, equivalent to 2.4\% (Table~\ref{tab:ch3-sample-sizes}, \S\ref{sec:ch3-data}). The robustness check re-runs the full pipeline under two alternative thresholds, $\alpha = 0.0001$ (more conservative; outliers removed = 14) and $\alpha = 0$ (no exclusion). The K-means partition at $k = 6$ is recomputed under each variant on the 31-variable standardised input space with the same seed and initialisation count of the deployed routine, and the adjusted Rand index is computed against the deployed partition over the intersection of the two samples. The resulting ARI values are 0.654 at $\alpha = 0.0001$ ($+30$ rows retained) and 0.551 at $\alpha = 0$ ($+44$ rows retained). The values lie below the 0.80 reference threshold for substantive partition robustness.

The lower-than-threshold ARI is qualified by a profile-level purity calculation that maps each variant cluster to the dominant deployed profile and computes the proportion of observations in the majority-profile column. The purity is 0.822 at $\alpha = 0.0001$ and 0.727 at $\alpha = 0$. The interpretation is that the six macro-profiles of the deployed typology survive both perturbations --- the Wealthy Digital, Connected Wealthy, Fragile and Asset Rich clusters retain identifiable variant-cluster homes under both alternatives --- while the boundary observations between contiguous profiles (Moderate, Social Decline) migrate across cluster assignments at the higher rate observed at $\alpha = 0$. The choice of $\alpha = 0.001$ is retained as the conventional threshold; the partition is read as robust at the macro-profile level under the stricter alternative and as moderately less robust under the no-exclusion alternative. The profile-level confusion of the deployed partition against the stricter $\alpha = 0.0001$ alternative is reported in \S\ref{sec:ch7-outlier-robustness} (Table~\ref{tab:ch7-outlier-robustness}).


\section{Distance metric: Gower versus Euclidean}
\label{sec:appx-sens-gower}

The deployed K-means routine for both countries operates on the Euclidean distance over the standardised variable space. The robustness check substitutes the Gower dissimilarity computed via \texttt{cluster::daisy} on the original (non-standardised) data with all binary variables declared as asymmetric, and runs partitioning-around-medoids at the deployed $k$ ($k = 5$ for Italy, $k = 6$ for Sweden). The adjusted Rand index between the Gower-PAM partition and the deployed K-means partition is 0.205 for Italy and 0.091 for Sweden. Both values are below the 0.70 reference threshold for partition convergence across distance choices.

The empirical disagreement between the two distance choices is read as substantively informative rather than as a failure of either algorithm. The Euclidean choice is retained on three grounds. The standardised input space carries an interpretation of cluster centroids as mean $z$-deviations from the population baseline, a quantity that is consumed directly by the ANOVA validation of \S\ref{sec:ch3-kmeans} and the fingerprint visualisations of \S\ref{sec:ch4-integrative-view}; the Gower dissimilarity, by normalising each variable to its own range, dissolves the cross-variable comparability that underpins the centroid reading. The K-means / Euclidean combination is the standard practice in the welfare-regime and typology literature on which the comparative reading of Chapter~\ref{ch:ch6} builds. The Gower / PAM combination produces a different partition because it weights the binary indicators differently --- the asymmetric-binary handling implies that two zero-coded observations are treated as not similar rather than as jointly absent, a choice that suppresses the role of co-absence in the digital and social blocks, where the deployed reading treats co-absence as a substantively informative configuration of low engagement. The two partitions therefore represent two different definitions of what counts as similar, not two competing estimates of an underlying ground truth.

Profile-level analysis of the Gower partition confirms that the disagreement is pervasive rather than confined to a single stratum: the global profile-level purity is 0.52 in Italy and 0.38 in Sweden, and no Italian or Swedish profile concentrates more than 70\% of its members on a single Gower-PAM cluster. The low purity is the profile-level signature of the same effect the adjusted Rand index quantifies --- the metric choice reorganises the partition broadly rather than isolating a robust subset of profiles --- consistent with the cross-method evidence (\S\ref{sec:ch7-cross-method}) that several cluster boundaries are geometrically thin in the standardised space.

A complementary sensitivity holds the Euclidean K-means specification but down-weights the binary indicators, to probe the amplification that within-country $z$-standardisation introduces on low-prevalence binary variables. The 11 binary indicators in the Swedish active set (9 in the Italian active set, after the zero-variance exclusion of \texttt{ac035d4} and \texttt{ac035d7}) are re-scaled by a factor $w < 1$, leaving the 20 rank-transformed non-binary variables unchanged, and the K-means partition at the deployed $k$ is re-derived under each weighting; the unit-weight case $w = 1$ reproduces the deployed partition exactly. At $w = 0.5$ the adjusted Rand index against the deployed partition is 0.370 for Italy and 0.335 for Sweden, and at $w = 1/\sqrt{2}$ it is 0.401 for Italy and 0.379 for Sweden. The binary indicators therefore carry substantial partition-determining weight: down-weighting them reassigns a non-trivial share of respondents. At the profile level the macro-typology is more stable than the adjusted Rand index alone implies --- the deployed profiles retain dominant homes at a purity of 0.64--0.67 (Italy) and 0.60--0.64 (Sweden) across the two weightings --- but the sensitivity is real and is not read as robustness. The deployed unit-weighted specification is retained on two substantive grounds the down-weighting clarifies rather than refutes: the binary indicators of digital and social engagement are constitutive of the profiles the typology is built to recover (the Wealthy Digital and Connected configurations are defined precisely by these items), so down-weighting them dissolves distinctions the analysis is designed to capture; and any down-weight magnitude is a free parameter with no measurement-theoretic anchor, whereas unit-weighted $z$-standardisation preserves the interpretation of cluster centroids as deviations on a single common scale across binary and non-binary indicators.

Taken together, the three components --- Gower-PAM as an alternative metric, the profile-level breakdown of the disagreement, and binary down-weighting as a parametric alternative --- converge on a single reading: the partition is sensitive to both the distance metric and the weighting of the binary indicators, and that sensitivity reinforces, rather than undermines, the scope condition stated in \S\ref{sec:ch7-synthesis} --- that the substantive findings are claims about the centroid configuration of the profiles and the relations among them, not about the assignment of marginal respondents, which is the level at which the metric and weighting choices exert their influence.


\section{Alternative factor extraction methods (Sweden)}
\label{sec:appx-sens-fa}

The deployed factor analysis for Sweden (\S\ref{sec:ch3-fa}) extracts six factors via maximum-likelihood factoring on the 31 active variables, with varimax rotation. An earlier principal-axis specification produced a Heywood case on \texttt{social\_integration} (loading $1.11$, communality $h^2 = 1.38$), which the maximum-likelihood specification resolves cleanly. Three alternative specifications are reported here to document the fit and clustering stability of the deployed choice. The first re-extracts the six factors via principal-axis factoring on the same 31 variables (the prior deployed specification); the second substitutes minimum-residual factoring; the third reverts to principal-axis factoring but on the 30 variables that exclude \texttt{social\_integration}.

The principal-axis specification on the 31 variables --- the prior deployed choice --- retains the Heywood case on \texttt{social\_integration}, with RMSEA = 0.033 and TLI = 0.897. The minimum-residual specification likewise retains the Heywood case, with RMSEA = 0.034 and TLI = 0.890, as expected from a principal-axis-type estimator. The deployed maximum-likelihood specification eliminates the Heywood case at effectively identical fit (RMSEA = 0.033, TLI = 0.897), indicating that the Heywood loading is specification-dependent rather than data-pathological. The 30-variable principal-axis specification, which drops \texttt{social\_integration}, also eliminates the Heywood case and returns the best overall fit of the alternatives (RMSEA = 0.029, TLI = 0.908, KMO = 0.830 versus 0.808 on the 31-variable input). The factor structure is therefore not unique: the choice of extraction method and the inclusion of \texttt{social\_integration} each shape whether the Heywood case appears, while every specification yields a defensible six-factor solution under standard fit criteria.

The implication for the deployed clustering is independent of the factor choice for the three 31-variable specifications. The K-means routine operates on the standardised variable space and is invariant to the factor extraction by construction --- the adjusted Rand index of the principal-axis and minimum-residual specifications against the deployed maximum-likelihood partition is unity, since the cluster input is identical. The 30-variable principal-axis specification does alter the K-means input space, and the resulting ARI is 0.518: the removal of \texttt{social\_integration} from the cluster input induces a substantively different partition over the Swedish sample. The deployed specification (ML, 31 variables, \texttt{social\_integration} retained) is therefore defensible on two grounds that the alternatives clarify rather than challenge --- the factor analysis is descriptive and does not propagate into the clustering, which is invariant to the extraction method, and the inclusion of \texttt{social\_integration} is required by the Isolation Paradox reading in Chapter~\ref{ch:ch4}, which the maximum-likelihood extraction now preserves without incurring the Heywood case of the principal-axis alternative.


\section{Italian 29-versus-31 variable asymmetry}
\label{sec:appx-sens-29-31}

The Italian analytical sample exhibits zero variance on the two activity-indicator variables \texttt{ac035d4} and \texttt{ac035d7}, which are consequently dropped from the active variable set, yielding 29 active variables against the 31 active variables of the Swedish branch (\S\ref{sec:ch3-fa}). The 29-versus-31 asymmetry is documented in the main text; this section reports the preemptive robustness check that quantifies the impact of the asymmetry on the deployed Italian partition.

The two zero-variance variables are re-included with a small Gaussian perturbation ($\sigma = 0.05$ on the standardised scale) drawn under the deployed seed; the perturbation magnitude is calibrated to introduce non-degenerate variance without injecting structure. The full pipeline is re-derived on the 31-variable input space, with principal-axis factor analysis at six factors and varimax rotation, K-means clustering at $k = 5$, and identical seed and initialisation specifications. The factor analysis converges cleanly with no Heywood case (RMSEA = 0.049, RMS = 0.027). The adjusted Rand index between the 31-variable variant partition and the deployed 29-variable partition is 0.977. The 29-versus-31 asymmetry is read as empirically benign: the deployed Italian partition is structurally identical, up to a negligible boundary perturbation, to the partition that would have been derived had the two zero-variance variables not been dropped.


\section{Summary}
\label{sec:appx-sens-summary}

Table~\ref{tab:appx-sens-summary} reports the four sensitivity analyses against the interpretive thresholds stated in the preceding sections. The 29-versus-31 variable asymmetry is closed empirically. The outlier-threshold sensitivity is read as a moderate concern at the boundary-observation level and as robust at the macro-profile level. The distance-metric and binary-weighting sensitivity is examined through a three-component analysis (Gower-PAM as an alternative metric, profile-level confusion, and binary down-weighting): the deployed partition is sensitive to both the distance metric and the weighting of the binary indicators, and the unit-weighted Euclidean specification is retained on the substantive grounds stated in \S\ref{sec:appx-sens-gower} rather than as a claim of invariance to these choices. The factor-extraction sensitivity is resolved in the deployed pipeline: Sweden adopts maximum-likelihood extraction (\S\ref{sec:appx-sens-fa}), which removes the Heywood case at unchanged fit and leaves the clustering invariant.

\begin{table}[ht]
\centering
\caption[Phase 2 sensitivity analyses summary]{Phase 2 sensitivity analyses --- summary against interpretive thresholds.}
\label{tab:appx-sens-summary}
\small
\setlength{\tabcolsep}{6pt}
\begin{tabular}{p{0.42\textwidth} p{0.36\textwidth} p{0.12\textwidth}}
\hline
\textbf{Analysis} & \textbf{Metric (vs.\ threshold)} & \textbf{Verdict} \\
\hline
Outlier exclusion threshold (Sweden, $\alpha = 0.0001$ / $\alpha = 0$) &
ARI = 0.654 / 0.551 vs.\ 0.80; profile-purity = 0.822 / 0.727 &
Concern (macro robust) \\
\\
Distance metric and binary weighting (Gower-PAM; profile-level confusion; binary down-weight $w = 0.5$ / $w = 1/\sqrt{2}$), Italy / Sweden &
Gower-PAM ARI 0.205 / 0.091 (purity 0.52 / 0.38). Down-weight ARI ($w{=}0.5$) 0.370 / 0.335, ($w{=}1/\sqrt{2}$) 0.401 / 0.379; profile purity 0.60--0.69. Sensitive to metric and binary weight &
Concern (defended, three-prong) \\
\\
Heywood case, Sweden FA (deployed PA $\to$ ML; alternatives PA-31v / minres-31v / PA-30v) &
Deployed ML: Heywood FALSE, RMSEA 0.033. Alternatives Heywood TRUE / TRUE / FALSE; RMSEA 0.033 / 0.034 / 0.029; ARI on clustering input 1.000 / 1.000 / 0.518 &
Resolved (ML deployed) \\
\\
Italian 29-vs-31 variable asymmetry (perturbed re-inclusion) &
ARI = 0.977 vs.\ 0.80; Heywood = FALSE &
Pass \\
\hline
\end{tabular}
\end{table}

\section{Unweighted versus post-stratified cluster shares}
\label{sec:appx-weighted-shares}

The segmentation is estimated on the unweighted analytical sample (\S\ref{sec:ch3-preprocessing}). To document the empirical magnitude of the SHARE survey design, a post-hoc projection re-weights the within-sample cluster fractions using the SHARE Wave~9 calibrated cross-sectional individual weight (\texttt{cciw\_w9}), computing each profile's weighted share as $\sum_{i \in C} w_i / \sum_i w_i$, under the additional assumption that the proxy-interview exclusion of the analytical sample does not bias the weighted-share estimate. Table~\ref{tab:ch3-weighted-shares} reports the Italian and Swedish cluster shares under the deployed unweighted specification and under this post-hoc re-weighting. The differences are uniformly within 3.2 percentage points across all eleven profiles, so the descriptive cluster shares of Chapters~\ref{ch:ch4} and~\ref{ch:ch5} are broadly stable under the weighting adjustment; the national-population projection of Figure~\ref{fig:ch4-shares} applies these weighted post-hoc shares to the 2024 Istat over-65 estimate.

\begin{table}[!htbp]
\centering
\caption[Unweighted vs weighted cluster shares]{Cluster shares under deployed unweighted specification and post-hoc weighted re-projection}
\label{tab:ch3-weighted-shares}
\small
\begin{tabular}{llrrr}
\toprule
\textbf{Country} & \textbf{Profile} & \textbf{Unweighted (\%)} & \textbf{Weighted (\%)} & \textbf{$\Delta$ pp} \\
\midrule
Italy  & Fragile Resigned    & 13.2 & 14.5 & $+1.2$ \\
Italy  & Fragile Depressed   & 26.7 & 27.0 & $+0.3$ \\
Italy  & Moderate Isolated   & 26.9 & 26.0 & $-0.9$ \\
Italy  & Traditional Social  &  7.7 &  8.6 & $+0.9$ \\
Italy  & Connected Active    & 25.4 & 23.9 & $-1.5$ \\
\midrule
Sweden & Fragile             & 11.5 & 13.0 & $+1.5$ \\
Sweden & Social Decline      &  8.3 &  9.0 & $+0.7$ \\
Sweden & Moderate            & 23.3 & 20.1 & $-3.2$ \\
Sweden & Asset Rich          & 19.1 & 20.8 & $+1.8$ \\
Sweden & Wealthy Digital     &  7.7 &  8.0 & $+0.3$ \\
Sweden & Connected Wealthy   & 30.0 & 29.0 & $-1.0$ \\
\bottomrule
\end{tabular}
\par\medskip
\begin{minipage}{\linewidth}
\footnotesize\textit{Note.} Unweighted shares are within-sample cluster fractions. Weighted shares are computed as $\sum_{i \in C} w_i / \sum_i w_i$ using the SHARE Wave~9 calibrated cross-sectional individual weight (\texttt{cciw\_w9}). The population projection in Figure~\ref{fig:ch4-shares} applies the weighted shares to the Istat 2024 Italian over-65 population estimate (14.18~million).
\end{minipage}
\end{table}
